% Nouveau EGA V (cote 25): archive cote 25, site page 27 (PDF page 28).
% typescript. Draft compared: none.
% AI-transcribed by Claude Opus 5.5 (claude-opus-5-5, adaptive thinking, effort high), Message Batches API, batch ega25b (msgbatch_01A7n5YuXubUPNK1zsepi8RG), 2026-11-13; original AI transcription, subsequently checked against both supplied scan parts; see PAGES_CHECKED.tsv and CHECK_DECISIONS.tsv.
\scanpage{25/27}
\mspage{\hand{23}}



7. \uline{Application à la construction de sections hyperplanes ou de multisections
de type précisé}.

7.1. Notons que si S=Spec(k), ou k est un corps \uline{infini}, alors tout ouvert
non vide de $\overset{\hand{\vee}}{P}$ contient un point rationnel sur k; donc avec les notations de
4.1. , si E (défini en termes d'une propriété constructible P) contient
le point générique \hand{$\eta$} , il cotient  un point rationnel sur k, donc
il existe une section hyperplane de X (définie sur k) ayant la propriété P.
D'autre part, S étant de nouveau quelconque, il est immédiat que pour tout
s\hand{$\in$}S et tout point $\xi_0$ de la fibre $\overset{\hand{\vee}}{P}_s$ rationnel sur k(s), il existe une section
\hand{$\xi$} de $\overset{\hand{\vee}}{P}$ sur un voisinage ouvert U de s qui passe par \hand{$\xi$}$_0$ . Si alors E est
encore  défini comme dans 4.1. en termes d'une propriété constructible
P, et si on dispose (p ex grâce au N°5) du fait que E est ouvert, alors
si \hand{$\xi$}$_0$\hand{$\in$} E,  la section \hand{$\xi$} est une section de E sur 
U, du moins si on rapetisse suffisamment U. Donc, on aura construit une
section hyperplane $Y_{\hand{\xi}}$ de X , au dessus d'un voisinage ouvert U de s, telle que

\placedmargin{27}
